% A fictional paper for the evals of the paper-website skill. Every number below is in records/; the scores are
% synthetic, drawn by a seeded generator (the same values as the starter's example results.json).
\documentclass{article}
\usepackage{amsmath}
\title{Kestrel: a pretrained encoder for multivariate sensor streams}
\author{Ana Example$^{1}$, Ben Sample$^{2}$, Chen Placeholder$^{1}$, Dora Fictive$^{1,\dagger}$\\
$^{1}$Institute of Example Studies, $^{2}$Sample University, $^{\dagger}$supervision}
\begin{document}
\maketitle

\begin{abstract}
Sensor streams from power grids, data centres, wearables, roads and machines are long, irregular and rarely labelled.
Models trained for one stream do not transfer to the next.
We present Kestrel, an encoder pretrained on 137 public datasets that turns any multivariate stream into one embedding per time step.
On five benchmark tasks Kestrel ranks first on four and second on one, with a rating of 1125 against 1025 for the strongest comparator.
\end{abstract}

\section{Introduction}
Every new sensor deployment starts without labels.
A model trained on one stream breaks when the sampling rate, the channel count or the noise changes.
Labelled anomalies are rare, so supervised models see a few dozen positives at most.
We ask whether one encoder, pretrained once on public streams, can serve every new stream without retraining.
Kestrel turns any multivariate stream into one embedding per time step, so a linear probe can read any downstream label from it.

\section{Method}
Kestrel cuts each channel into patches of 24 steps and embeds each patch with a shared tokenizer.
A cross-channel attention layer mixes the channels at every patch, so the encoder accepts any number of channels.
During pretraining, 35\% of the patches are masked and the encoder predicts their embeddings from the rest.
A masking curriculum grows the masked span from one patch to six over training.
Kestrel was pretrained on 137 public datasets with 2.4 billion time steps from 7 domains.
The model has 31.6 million parameters in 10 layers.
This is the first checkpoint; larger models trained on more domains will follow.

\section{Results}
We compare Kestrel with four methods on five tasks: grid load forecasting, server log anomaly detection, gait phase classification, traffic incident detection and bearing wear grading.
Every method is read with the same linear probe on frozen embeddings, and scored by macro F1.
\begin{itemize}
\item Grid load: forecast the next hour of load from the previous day, scored as a three-class trend.
\item Server logs: flag the minutes that operators marked as incidents.
\item Gait: label each step phase from a wrist sensor.
\item Traffic: detect incidents from loop detector counts.
\item Bearing wear: grade the wear of a bearing from its vibration, in three classes.
\end{itemize}
Kestrel ranks first on four of the five tasks and second on gait.
Across tasks, Kestrel has a rating of 1125 against 1025 for the strongest comparator, Method B, where the baseline is fixed at 1000.
Kestrel is rated above every other method in a two-sided sign test over all task and seed pairs ($p = 0.002$).
On traffic alone the difference to Method B is weaker ($p = 0.04$; 95\% confidence interval of the F1 difference 0.003 to 0.047).

\begin{table}
\caption{Macro F1 per task on held-out streams; mean rank over the five tasks.}
\begin{tabular}{lrrrrrr}
method & grid load & server logs & gait & traffic & bearing wear & mean rank\\
Kestrel & .587 & .773 & .626 & .613 & .709 & 1.20\\
Method B & .381 & .673 & .767 & .588 & .508 & 3.00\\
Method C & .523 & .582 & .517 & .411 & .646 & 3.20\\
Method D & .511 & .654 & .451 & .374 & .532 & 4.00\\
Baseline & .473 & .379 & .557 & .544 & .549 & 3.60\\
\end{tabular}
\end{table}

\section{Ablations}
Removing any part of Kestrel lowers macro F1 on at least one task at one of the two model sizes.
Without pretraining, F1 falls by up to .157; without the masking curriculum, by up to .094.
The cross-channel attention matters least: removing it costs at most .056.

\section{Limitations}
Kestrel saw only public datasets in pretraining, and seven domains do not cover every sensor.
Streams sampled faster than 2 kHz were not part of pretraining.
The five tasks use linear probes; fine-tuning was not evaluated.

\section{Conclusion}
One pretrained encoder can read many kinds of sensor streams without being trained again.

\end{document}
